% =====================================================================
% RELATÓRIO 6 — POO com C++: seção de exercícios (versão de 2026.2)
% Substitui a seção "Exercícios Práticos do Módulo" do capítulo 1. O texto
% segue os slides enviados aos alunos (Relatório 06/poo_cpp.pdf).
%
% As caixas usam o ambiente exercicioCpp, com o mesmo visual do ambiente
% exercicio. Esse último só é criado no capítulo do Relatório 7, sem
% proteção contra redefinição; por isso este arquivo não o define nem o usa.
% Se o capítulo 1 já tiver um ambiente próprio para as caixas, basta
% trocar \begin{exercicioCpp}{...} por ele.
% =====================================================================
\definecolor{azulRelatorio}{rgb}{0.1,0.3,0.8}

\ifcsname exercicioCpp\endcsname\else
\newtcolorbox{exercicioCpp}[1]{
    colback=white,
    colframe=azulRelatorio,
    coltitle=white,
    fonttitle=\normalfont,
    boxrule=0.8pt,
    arc=1.5mm,
    title={#1}
}
\fi

\section{Exercícios Práticos do Módulo}

\begin{exercicioCpp}{Exercício 1: Classes, Objetos e Métodos (Duelo de Bandas)}
Em um grande festival de música com um duelo de bandas, cada banda possui um nome, uma
quantidade de integrantes, uma potência de som e um nível de energia da plateia. Implemente:

\begin{itemize}
    \item \texttt{Banda}: atributos \texttt{nome} (\texttt{string}), \texttt{integrantes}
    (\texttt{int}), \texttt{potenciaSom} (\texttt{float}) e \texttt{energia} (\texttt{int}); e o
    método \texttt{duelar()}, que recebe como parâmetro um objeto da classe \texttt{Banda} (a
    rival), exibe uma mensagem confirmando a apresentação e subtrai o valor da
    \texttt{potenciaSom} da \texttt{energia} da banda rival. Para que a energia da rival seja
    alterada de fato, receba-a por referência: \texttt{duelar(Banda \&rival)}.
\end{itemize}

No programa principal, instancie dois objetos da classe \texttt{Banda} e atribua valores
personalizados aos atributos de cada um. Execute o método \texttt{duelar()}, definindo qual
banda será a desafiante e qual será a rival. Ao final, exiba o status atualizado das duas
bandas, para ver como ficou a energia após o confronto.
\end{exercicioCpp}

\begin{exercicioCpp}{Exercício 2: Encapsulamento e Modificadores de Acesso (Links Sociais)}
No jogo \textit{Persona}, a relação do protagonista com os seus aliados (os \textbf{Links
Sociais}) evolui à medida que eles passam tempo juntos, alterando os seus níveis de afinidade
de forma controlada. Implemente:

\begin{itemize}
    \item \texttt{LinkSocial}: atributos privados \texttt{nome} (\texttt{string}, o nome do
    personagem), \texttt{arcana} (\texttt{string}, a arcana que representa o Link entre o
    protagonista e o personagem) e \texttt{rank} (\texttt{int}); métodos públicos de acesso
    (\textit{getters}) e de modificação (\textit{setters}) para cada atributo; e o método
    \texttt{subirRank()}, que apenas incrementa em 1 o valor do atributo \texttt{rank}.
\end{itemize}

No programa principal, instancie um objeto da classe \texttt{LinkSocial} e use os
\textit{setters} para definir o nome, a arcana e o rank inicial (por exemplo, 1) do personagem.
Execute o método \texttt{subirRank()} e exiba os dados do Link Social usando os
\textit{getters}, para confirmar o aumento do rank.
\end{exercicioCpp}

\begin{exercicioCpp}{Exercício 3: Herança e Sobrescrita (Membros do Inatel)}
No Inatel, durante eventos acadêmicos e feiras tecnológicas como a FETIN, diferentes pessoas
desempenham papéis essenciais na comunidade. Todos compartilham características de um membro
da instituição, mas alunos e professores possuem atribuições e dados específicos dentro do
campus. Implemente:

\begin{itemize}
    \item \texttt{MembroInatel} (classe base): atributo \texttt{nome} (\texttt{string}) e método
    \texttt{seApresentar()}, que exibe ``Sou um membro da comunidade Inatel: [nome].''

    \item \texttt{Aluno}: herda de \texttt{MembroInatel}, tem o atributo \texttt{curso}
    (\texttt{string}) e sobrescreve \texttt{seApresentar()} para exibir ``Meu nome é [nome] e
    estudo no curso de [curso].''

    \item \texttt{Professor}: herda de \texttt{MembroInatel}, tem o atributo
    \texttt{disciplina} (\texttt{string}) e sobrescreve \texttt{seApresentar()} para exibir
    ``Meu nome é [nome] e leciono a disciplina de [disciplina].''
\end{itemize}

No programa principal, instancie um objeto da classe \texttt{Aluno} e um da classe
\texttt{Professor} e atribua valores aos atributos de cada um. Chame o método
\texttt{seApresentar()} de cada objeto, para demonstrar como a classe filha reutilizou e
personalizou o comportamento da classe base.
\end{exercicioCpp}

\begin{exercicioCpp}{Exercício 4: Polimorfismo em Coleções (Ocupações da Comarca)}
Na \textbf{Comarca}, a terra natal dos Hobbits, a vida é tranquila e cheia de tradições
cotidianas. Embora todos compartilhem o amor pela paz e pela boa comida, cada Hobbit possui uma
ocupação diferente para manter a comunidade funcionando. Implemente:

\begin{itemize}
    \item \texttt{Hobbit} (classe base): atributo \texttt{nome} (\texttt{string}) e método
    virtual \texttt{fazerAtividade()}, que imprime ``O hobbit [nome] está aproveitando um dia
    tranquilo na Comarca.''

    \item \texttt{Jardineiro}: herda de \texttt{Hobbit} e sobrescreve
    \texttt{fazerAtividade()} para exibir ``O jardineiro [nome] está cuidando das flores e
    plantas ao redor das tocas!''

    \item \texttt{Cozinheiro}: herda de \texttt{Hobbit} e sobrescreve
    \texttt{fazerAtividade()} para exibir ``O cozinheiro [nome] está preparando o segundo café
    da manhã para os convidados!''

    \item \texttt{Fazendeiro}: herda de \texttt{Hobbit} e sobrescreve
    \texttt{fazerAtividade()} para exibir ``O fazendeiro [nome] está colhendo vegetais e
    hortaliças em suas terras!''
\end{itemize}

No programa principal, crie um vetor (ou lista) para armazenar os objetos do tipo
\texttt{Hobbit}, usando ponteiros da classe base (\texttt{std::vector<Hobbit*>}). Crie uma
instância de cada profissão (\texttt{Jardineiro}, \texttt{Cozinheiro} e \texttt{Fazendeiro}),
com os seus respectivos nomes, e adicione-as ao vetor. Percorra o vetor e, para cada elemento,
chame o método \texttt{fazerAtividade()}, para demonstrar o comportamento polimórfico.
\end{exercicioCpp}
